% =====================================================================
%  How-To unificado — Acceso Remoto Seguro con WireGuard sobre the-store
%  Documento autocontenido. Compilar con:  pdflatex how-to.tex  (x2)
%  o:  latexmk -pdf how-to.tex
% =====================================================================
\documentclass[11pt,a4paper]{article}

\usepackage[utf8]{inputenc}
\usepackage[T1]{fontenc}
\usepackage[spanish,es-noshorthands,es-tabla]{babel}
\usepackage[a4paper,margin=2.4cm]{geometry}
\usepackage{lmodern}
\usepackage{microtype}
\usepackage{xcolor}
\usepackage{graphicx}
\usepackage{booktabs}
\usepackage{tabularx}
\usepackage{array}
\usepackage{enumitem}
\usepackage{listings}
\usepackage{fancyhdr}
\usepackage{titlesec}
\usepackage[hidelinks,colorlinks=true,linkcolor=azulvpn,urlcolor=azulvpn,citecolor=azulvpn]{hyperref}

% ---- Paleta ----
\definecolor{azulvpn}{HTML}{1F4E79}
\definecolor{grisfondo}{HTML}{F4F6F8}
\definecolor{grisborde}{HTML}{D0D7DE}
\definecolor{verdeok}{HTML}{1A7F37}
\definecolor{rojoko}{HTML}{B42318}
\definecolor{comentario}{HTML}{6A737D}
\definecolor{keyword}{HTML}{8250DF}
\definecolor{string}{HTML}{0A3069}

% ---- Encabezado / pie ----
\pagestyle{fancy}
\fancyhf{}
\renewcommand{\headrulewidth}{0.4pt}
\renewcommand{\footrulewidth}{0pt}
\fancyhead[L]{\small\textcolor{azulvpn}{TPE Redes — WireGuard + Kubernetes}}
\fancyhead[R]{\small\leftmark}
\fancyfoot[C]{\small\thepage}

% ---- Títulos ----
\titleformat{\section}{\Large\bfseries\color{azulvpn}}{\thesection}{0.6em}{}
\titleformat{\subsection}{\large\bfseries\color{azulvpn!85!black}}{\thesubsection}{0.6em}{}
\titleformat{\subsubsection}{\normalsize\bfseries}{\thesubsubsection}{0.6em}{}

% ---- Estilo de código ----
\lstdefinestyle{shell}{
  basicstyle=\ttfamily\small,
  backgroundcolor=\color{grisfondo},
  frame=single,
  framesep=6pt,
  rulecolor=\color{grisborde},
  breaklines=true,
  breakatwhitespace=false,
  postbreak=\mbox{\textcolor{comentario}{$\hookrightarrow$}\space},
  showstringspaces=false,
  columns=fullflexible,
  keepspaces=true,
  upquote=true,
  commentstyle=\color{comentario}\itshape,
  morecomment=[l]{\#},
}
\lstset{style=shell}

\newcommand{\cmd}[1]{\texttt{#1}}
\newcommand{\si}{\textcolor{verdeok}{\textbf{\checkmark}}}
\newcommand{\no}{\textcolor{rojoko}{\textbf{\ding{55}}}}
\usepackage{pifont}
\renewcommand{\si}{\textcolor{verdeok}{\ding{51}}}
\renewcommand{\no}{\textcolor{rojoko}{\ding{55}}}

\setlist{itemsep=2pt,topsep=3pt}

% ---- Índice: más ancho para el número de sección (babel agrega un punto,
%      y "10." no entra en los 1.5em por defecto de article) ----
\makeatletter
\renewcommand*\l@section[2]{%
  \ifnum \c@tocdepth >\z@
    \addpenalty\@secpenalty
    \addvspace{1.0em \@plus\p@}%
    \setlength\@tempdima{2.3em}%
    \begingroup
      \parindent \z@ \rightskip \@pnumwidth
      \parfillskip -\@pnumwidth
      \leavevmode \bfseries
      \advance\leftskip\@tempdima
      \hskip -\leftskip
      #1\nobreak\hfil
      \nobreak\hb@xt@\@pnumwidth{\hss #2}\par
    \endgroup
  \fi}
\makeatother

% =====================================================================
\begin{document}

% --------------------------- CARÁTULA --------------------------------
\begin{titlepage}
\thispagestyle{empty}
\centering
\vspace*{2.2cm}

{\color{azulvpn}\rule{\linewidth}{1.2pt}}\\[0.5cm]
{\Huge\bfseries\color{azulvpn} Acceso Remoto Seguro\\[0.25cm] con WireGuard sobre \textit{the-store}}\\[0.4cm]
{\large Despliegue, pruebas y alcance acotado de la solución}\\[0.3cm]
{\color{azulvpn}\rule{\linewidth}{1.2pt}}\\[1.4cm]

{\large\bfseries Trabajo Práctico Especial — Redes de Información}\\[0.4cm]
{\large WireGuard (Client-To-Site / Site-To-Site) + Kubernetes NetworkPolicies}\\[2.2cm]

\begin{tabular}{rl}
\textbf{Plataforma:} & Kind (cluster de una sola PC) \\[4pt]
\textbf{Aplicación:}  & the-store (5 microservicios) \\[4pt]
\textbf{Documento:}   & How-To unificado + problemática \\
\end{tabular}

\vfill
{\large\today}
\end{titlepage}

% --------------------------- RESUMEN ---------------------------------
\phantomsection
\section*{Resumen}
\addcontentsline{toc}{section}{Resumen}
Esta guía documenta, de punta a punta, una solución de \textbf{acceso remoto
seguro} al cluster de microservicios \textit{the-store}: cómo desplegarla, cómo
probarla y \textbf{qué problemas resuelve y cuáles no}. La
solución combina \textbf{WireGuard} (túneles cifrados Client-To-Site y
Site-To-Site) con \textbf{Kubernetes NetworkPolicies} (segmentación de red).

El documento está dividido en dos partes. La \textbf{Parte~I} la
problemática. La \textbf{Parte~II} es el how-to operativo:
despliegue paso a paso y demostraciones reproducibles.

\newpage
\tableofcontents
\newpage

% =====================================================================
\phantomsection
\part*{Parte I — Problemática}
\addcontentsline{toc}{part}{Parte I — Problemática}
% =====================================================================

\section{Contexto real del deploy}
\textit{The Store} corre sobre un cluster \textbf{Kind} (una sola máquina) con 5
microservicios: UI, Catalog, Cart, Orders y Checkout. El estado verificado sobre
el repositorio (\cmd{dist/kubernetes.yaml}, \cmd{local.sh}) es:

\begin{itemize}
  \item \textbf{Todos los Services son \cmd{ClusterIP}}: no hay NodePort ni
        LoadBalancer en la aplicación.
  \item \textbf{El único punto de entrada es el Ingress NGINX} publicado en
        \cmd{localhost:80} (sólo se mapean 80/443 del nodo al host).
  \item \textbf{No existen NetworkPolicies de fábrica}: cualquier pod alcanza a
        cualquier pod.
  \item \textbf{Persistencia \cmd{in-memory}} en los 5 servicios: no hay
        Redis/Postgres/RabbitMQ desplegados sobre la red.
\end{itemize}

\noindent\textbf{Consecuencia central:} fuera del Ingress en \cmd{localhost:80},
el cluster \textbf{no es alcanzable desde fuera del nodo de ninguna forma
segura}. No hay canal para operarlo remotamente ni para que una red externa
consuma sus servicios.

\section{Problema}
Acotamos el trabajo a \textbf{dos problemas de acceso remoto} y \textbf{uno de
segmentación} que los habilita de forma segura.

\subsection*{Problema A — No hay acceso administrativo remoto seguro}
Un administrador o desarrollador \textbf{fuera de la red del nodo} no puede operar
el cluster (\cmd{kubectl} contra el API server, ni alcanzar servicios internos)
sin exponer el API server o los servicios en texto plano.
\;$\rightarrow$\; \textbf{Resuelto con WireGuard Client-To-Site.}

\subsection*{Problema B — Una red externa no puede consumir los servicios}
Una \textbf{red corporativa completa} necesita consumir los servicios del cluster
\textbf{sin instalar un cliente VPN en cada equipo} y sin abrir el cluster a la
red. \;$\rightarrow$\; \textbf{Resuelto con WireGuard Site-To-Site} (túnel
permanente entre el gateway corporativo y el del cluster; los equipos finales
sólo necesitan una ruta estática).

\subsection*{Problema C — Sin segmentación, el acceso remoto es acceso total}
Una vez que el túnel deja entrar tráfico a la red del cluster, \textbf{cualquier
origen alcanzaría cualquier pod} (movimiento lateral). \;$\rightarrow$\;
\textbf{Resuelto con NetworkPolicies} (default-deny + reglas mínimas por
servicio). Es la pieza de defensa en profundidad que hace que el túnel habilite
\emph{sólo} lo necesario.

\subsection*{Qué provee cada tecnología (sin sobre-prometer)}
\renewcommand{\arraystretch}{1.3}
\begin{tabularx}{\linewidth}{@{}>{\bfseries}p{2.8cm} X@{}}
\toprule
Tecnología & Qué aporta exactamente \\
\midrule
WireGuard & Canal de red \textbf{cifrado (ChaCha20-Poly1305) y autenticado por
clave pública} entre peers, para los tramos \emph{cliente$\leftrightarrow$gateway}
y \emph{red corporativa$\leftrightarrow$gateway}. \cmd{AllowedIPs} define qué
subredes cruzan el túnel (mínimo privilegio de ruteo). \\
NetworkPolicies & Segmentación \textbf{a nivel de red del cluster}: default-deny y
reglas explícitas de qué servicio habla con qué servicio. Bloquea movimiento
lateral. \\
\bottomrule
\end{tabularx}

\section{POC y casos de uso}
\textbf{Incluye (demostrable, todo en una PC con Kind):} gateway VPN del lado del
cluster; Client-To-Site (admin $\rightarrow$ \cmd{kubectl} + servicios);
Site-To-Site (red corporativa $\rightarrow$ servicios sin cliente VPN en el
equipo final); NetworkPolicies (default-deny + reglas mínimas); validación de
tráfico cifrado y rotación de claves.

\begin{description}[leftmargin=1.4cm,style=nextline]
  \item[CU-1] \textbf{Admin opera el cluster remotamente (Client-To-Site):}
        levanta su túnel, hace \cmd{kubectl} contra el API server y alcanza
        servicios internos; sólo llega a las subredes de sus \cmd{AllowedIPs}.
  \item[CU-2] \textbf{Empleado consume servicios sin cliente VPN (Site-To-Site):}
        desde la red corporativa, el tráfico hacia un servicio cruza el túnel de
        forma transparente. El equipo final \textbf{no tiene WireGuard}.
  \item[CU-3] \textbf{Segmentación efectiva (NetworkPolicies):} un origen que
        entró por el túnel \textbf{no} puede hablar con un servicio no permitido
        (p.\,ej. alcanzar Orders salteándose a UI).
  \item[CU-4] \textbf{Rotación de claves ante compromiso:} se rota el par de un
        peer; las claves anteriores quedan invalidadas y el túnel se restablece.
\end{description}

% =====================================================================
\newpage
\phantomsection
\part*{Parte II — How-To: despliegue y pruebas}
\addcontentsline{toc}{part}{Parte II — How-To: despliegue y pruebas}
% =====================================================================

\section{Qué se demuestra}
\renewcommand{\arraystretch}{1.3}
\begin{tabularx}{\linewidth}{@{}>{\bfseries}p{3.2cm} X@{}}
\toprule
Escenario & Qué prueba \\
\midrule
Client-To-Site & Un admin remoto levanta un túnel y opera el cluster con
\cmd{kubectl} (API server \cmd{10.96.0.1}) y accede a servicios internos. Todo
cifrado. \\
Site-To-Site & Una red corporativa llega a los servicios \textbf{sin WireGuard en
el equipo final} (sólo ruta estática). \\
NetworkPolicies & default-deny + reglas mínimas: el túnel sólo alcanza lo
permitido; el movimiento lateral queda bloqueado. \\
Cifrado & En el cable sólo se ve UDP cifrado (ChaCha20-Poly1305); el HTTP viaja
en claro únicamente dentro del túnel. \\
Rotación de claves & Se rota un par en caliente; la clave anterior queda
invalidada. \\
\bottomrule
\end{tabularx}

\section{Prerequisitos}
\begin{itemize}
  \item \textbf{Docker}, \textbf{Kind} y \textbf{kubectl} instalados (ver
        \cmd{the-store-main/README.md}).
  \item Módulo \cmd{wireguard} del kernel (modo primario). Si no está disponible,
        la imagen cae a \textbf{userspace (\cmd{wireguard-go})} automáticamente.
        El módulo se carga en el kernel \textbf{del host} (los pods comparten ese
        kernel; no tienen kernel propio). Para forzar la carga:
        \cmd{sudo modprobe wireguard}.
\end{itemize}

\section{Arranque rápido (TL;DR)}
\begin{lstlisting}
# 1. App de microservicios (cluster Kind + the-store)
cd the-store-main && ./local.sh create-cluster --skip-tests && cd ..

# 2. Toda la solucion de acceso remoto, en orden, con verificacion
vpn/scripts/up.sh
\end{lstlisting}

\noindent\cmd{up.sh} corre, en orden: \cmd{00-gen-keys} $\rightarrow$
\cmd{10-deploy-gateway} $\rightarrow$ \cmd{20-up-admin-client} $\rightarrow$
\cmd{25-up-external-pc} $\rightarrow$ \cmd{30-up-site-to-site} $\rightarrow$
\cmd{40-apply-netpol} $\rightarrow$ \cmd{90-verify}. Al final debe imprimir
\textbf{\textcolor{verdeok}{RESULTADO: 10 PASS, 0 FAIL}} (con todos los escenarios
levantados).

\section{Paso a paso}

\subsection{Generar claves}
\begin{lstlisting}
vpn/scripts/00-gen-keys.sh
\end{lstlisting}
Crea 4 pares en \cmd{vpn/keys/} (gitignored): \cmd{gw-c2s}, \cmd{gw-s2s},
\cmd{admin}, \cmd{corp}. Cada privada nunca sale de su lado; sólo las públicas se
intercambian.

\subsection{Desplegar el gateway WireGuard}
\begin{lstlisting}
vpn/scripts/10-deploy-gateway.sh
kubectl get pod -n vpn-system
kubectl exec -n vpn-system deploy/wg-gateway -c wg-gateway -- wg show
\end{lstlisting}
Crea el namespace \cmd{vpn-system}, el Secret con \cmd{wg0.conf}/\cmd{wg1.conf} y
el Deployment + Service \textbf{NodePort UDP} (31820 Site-To-Site, 31821
Client-To-Site). El pod levanta \cmd{wg0} y \cmd{wg1}, habilita \cmd{ip\_forward}
(initContainer) y NATea hacia los CIDRs del cluster.

\subsection{Cliente admin (Client-To-Site)}
\begin{lstlisting}
vpn/scripts/20-up-admin-client.sh
docker exec admin wg show
docker exec admin kubectl --kubeconfig /root/admin.kubeconfig get pods -n the-store
\end{lstlisting}
El \cmd{kubectl} sale al API server \textbf{\cmd{10.96.0.1} por el túnel} (la ruta
a \cmd{10.96.0.0/16} sólo existe sobre \cmd{wg0}).

\subsection{Red corporativa (Site-To-Site)}
\begin{lstlisting}
vpn/scripts/30-up-site-to-site.sh
docker exec corp-pc wg show          # vacio: NO tiene WireGuard
CAT=$(kubectl get svc catalog -n the-store -o jsonpath='{.spec.clusterIP}')
docker exec corp-pc curl -s -o /dev/null -w '%{http_code}\n' http://$CAT/health   # 200
\end{lstlisting}

\subsection{Segmentación (NetworkPolicies)}
\begin{lstlisting}
vpn/scripts/40-apply-netpol.sh
kubectl get networkpolicy -n the-store
\end{lstlisting}

\subsection{Verificación end-to-end}
\begin{lstlisting}
vpn/scripts/90-verify.sh     # imprime PASS/FAIL por chequeo
\end{lstlisting}

\subsection{Rotación de claves (ante compromiso de un peer)}
\begin{lstlisting}
vpn/scripts/rotate-keys.sh admin    # o: rotate-keys.sh corp
\end{lstlisting}
Rota el par de claves del peer indicado \textbf{sin reiniciar el gateway}:
\begin{enumerate}
  \item Genera un par nuevo para el peer (\cmd{wg genkey} / \cmd{wg pubkey}).
  \item Actualiza el peer en el gateway \textbf{en caliente} con \cmd{wg set}:
        quita la clave pública vieja y agrega la nueva (la anterior queda
        \textbf{invalidada de inmediato}).
  \item Persiste el cambio en el Secret del gateway (sobrevive a un reinicio
        del pod).
  \item Reconfigura el cliente (\cmd{admin} o \cmd{corp-gateway}) con la clave
        nueva y reconecta el túnel.
\end{enumerate}
Al final el script verifica que la clave vieja \textbf{ya no figura} en el
gateway y que la nueva está activa (CU-4). Para comprobarlo a mano:
\begin{lstlisting}
kubectl exec -n vpn-system deploy/wg-gateway -c wg-gateway -- wg show wg0
docker exec admin wg show            # handshake reciente con la clave nueva
\end{lstlisting}

\section{Cómo funciona por dentro}
\begin{itemize}
  \item \textbf{Gateway = pod} en \cmd{vpn-system}, modo \textbf{kernel} (sólo
        \cmd{NET\_ADMIN}); un initContainer privilegiado habilita
        \cmd{ip\_forward} (en un pod \cmd{/proc/sys/net} es read-only). Fallback
        userspace (\cmd{wireguard-go}) si no hay módulo.
  \item \textbf{Exposición = NodePort UDP}: los clientes (containers en la red
        \cmd{kind}) alcanzan \cmd{<nodeIP>:31820/31821}.
  \item \textbf{AllowedIPs} = mínimo privilegio de ruteo y ACL de origen: definen
        qué subredes cruzan el túnel por cada peer.
  \item \textbf{NAT por destino} en el gateway: todo lo que se reenvía a
        \cmd{10.96.0.0/16} / \cmd{10.244.0.0/16} sale con la IP del pod, así el
        cluster sabe responder (vale para admin, túnel y la LAN corporativa).
  \item \textbf{NetworkPolicies} las aplica el CNI de Kind (kindnet): default-deny
        ingress + una regla por servicio.
\end{itemize}
\noindent Todos los valores de red están en \cmd{vpn/scripts/lib.sh}.

\subsection*{Mapa de rangos de red}
\renewcommand{\arraystretch}{1.3}
\begin{tabularx}{\linewidth}{@{}>{\ttfamily}p{3.2cm} X@{}}
\toprule
\normalfont\textbf{CIDR} & \textbf{Rol} \\
\midrule
10.200.0.0/24 & Subred del túnel Client-To-Site (IPs virtuales de WireGuard;
gateway \cmd{.1}, admin \cmd{.2}). \\
10.100.0.0/30 & Subred del túnel Site-To-Site (gateway corp $\leftrightarrow$
gateway cluster). \\
10.96.0.0/16  & Service CIDR del cluster (ClusterIPs; API server \cmd{10.96.0.1}). \\
10.244.0.0/16 & Pod CIDR del cluster (IPs reales de los pods). \\
172.21.0.0/16 & Red Docker \cmd{kind} (nodo y contenedores admin/corp/external). \\
172.30.0.0/24 & LAN corporativa simulada (\cmd{corp-pc}). \\
\bottomrule
\end{tabularx}

\section{Teardown}
\begin{lstlisting}
vpn/scripts/99-teardown.sh          # quita VPN/netpols, deja la app intacta
vpn/scripts/99-teardown.sh --all    # ademas borra las claves
\end{lstlisting}

% =====================================================================

\end{document}
